% Consolidated findings for the Moment-SoS-MINLP project.
%
% NOTE: this file is deliberately SEPARATE from notes/notes.tex, which is a one-way sync from Sam
% Chevalier's Overleaf repository (SamChevalier/Moment-SoS-MINLP-Notes) and whose contents can be
% overwritten by the next sync. Edit this file freely; coordinate with Sam before merging any of it
% into notes.tex or Overleaf.
%
% Compile:  pdflatex findings.tex
\documentclass[11pt]{article}
\usepackage[margin=1in]{geometry}
\usepackage{amsmath,amssymb}
\usepackage{booktabs}
\usepackage{longtable}
\usepackage{xcolor}
\usepackage[hidelinks]{hyperref}
\usepackage{siunitx}
\sisetup{group-separator={,},group-minimum-digits=4,detect-weight=true}

\newcommand{\code}[1]{\texttt{\small #1}}

\title{Moment/SOS relaxations for AC optimal transmission switching:\\consolidated findings}
\author{Daniel Molzahn \and Sam Chevalier}
\date{21 September 2026}

\begin{document}
\maketitle

\begin{abstract}
We apply sparse moment/sum-of-squares (Lasserre) relaxations to AC optimal transmission switching
(AC-OTS) on the 18 PGLib-OPF ``api'' test cases of Taheri and Molzahn, with every branch switchable.
Two threads are pursued: certified lower bounds from per-clique mixed-order relaxations, and correlated
randomized rounding of the switching binaries from the relaxation's pseudo-moments. We obtain certified
gaps below \SI{1}{\percent} on eight of the eighteen cases and below \SI{0.5}{\percent} on five, and we
find solutions that improve on the reference values for four cases. We report three findings that we
believe matter beyond this benchmark: (i) certification quality, not relaxation strength, drove an
earlier incorrect ranking of big-M versus no-big-M formulations; (ii) a network preprocessing step
(low-impedance bus merging) silently destroyed the validity of the bounds, producing a ``lower bound''
\SI{72}{\percent} above a known feasible solution, and removing it changes the certified gap by
between $+0.04$ and $-0.07$ percentage points --- it is free, and on the largest case an improvement; and (iii) the collapse of randomized rounding above 118 buses is
caused by the \emph{number} of lines opened simultaneously, not by network islanding, and capping it
raises the AC-feasible sample rate on \code{300-ieee} from \SI{0}{\percent} to \SI{91}{\percent} and
yields two new best-known solutions, on \code{179-goc} and \code{500-goc}.
\end{abstract}

\section{Scope}

The problem is AC-OTS written as a polynomial optimization problem in rectangular voltage coordinates
(see \code{docs/formulations.md}). Binaries $z_\ell \in \{0,1\}$ select branch status. The benchmark is
the set of 18 PGLib-OPF v21.07 ``api'' (congested) cases used by Taheri and Molzahn, with every branch
switchable; their AC-OPF, O-DC-OTS and AC-OTS (Juniper) values serve as reference points.

Two research threads:
\begin{enumerate}
  \item \textbf{Certified lower bounds} from sparse, per-clique mixed-order moment relaxations.
  \item \textbf{Correlated randomized rounding} of the binaries from the relaxation's pseudo-moments,
        with continuous recovery by Ipopt.
\end{enumerate}

\section{Formulation and method}

\paragraph{Binaries by basis reduction.} We impose $z \in \{0,1\}$ by substituting $z^k \to z$ in the
monomial basis rather than adding the equality $z^2 = z$. Keeping it as an equality makes the SDP
feasible set lower-dimensional, Slater's condition fails, and the order-2 interior-point solve degrades
badly. This is a practical, not cosmetic, distinction.

\paragraph{Switching is degree three.} The flow equations become $p_\ell = z_\ell \, P_\ell(V)$, which is
cubic. Order-1 (Shor) relaxations cannot represent this, and the resulting bounds collapse to
uselessness (\num{78116} on the 73-bus case, $0$ on the 118-bus case, \num{20079} on the 200-bus case).

\paragraph{Order policy rather than a global order.} Only cliques containing a binary are raised to
order 2 (\code{binary\_clique\_order}); all others remain at order 1. In practice 40--48\% of cliques are
order 2 and the largest PSD block has size 54--75, which is what makes the degree-3 constraints
affordable at these network sizes.

\paragraph{Variants.} \code{mixed} (big-M), \code{mixed\_nobigM}, \code{adj16} (adjacent-clique merging,
maximum size 16), \code{adj16\_pairs} (additional order-2 blocks on binary pairs), and, from
Experiment~4, \code{base} and \code{conn} (radial fixings plus connectivity cuts).

\paragraph{Certified bounds from inexact solves.} MOSEK terminates at \code{SLOW\_PROGRESS} on these
problems, so its dual objective is \emph{not} a valid bound. Following Oustry et al.~\cite{oustry2022}, we
compute a bound
valid at \emph{any} dual point $\theta$:
\[
  F(\theta) \;=\; \theta_t \;-\; \sum_\alpha |r_\alpha(\theta)| \max|x^\alpha|
              \;+\; \sum_k \rho_k \min\bigl(\lambda_{\min}(A_k(\theta)),\,0\bigr),
\]
where all correction terms are non-positive. Four variants are implemented (box; box with tight trace
bounds; implicit-Gram; per-block hybrid) plus an L-BFGS-optimized certificate on a smoothed surrogate.
Because $F$ is valid for every $\theta$, a badly converged solve still yields a rigorous bound.

\section{Main results}

Costs in \$/h. ``Best known'' is the minimum of our best found solution and the paper's AC-OPF, O-DC-OTS
and AC-OTS values. Gap is (best known $-$ certified bound)\,/\,best known. Every bound below is checked
against the best known feasible cost; none violates it.

\begin{longtable}{lrrrrrrl}
\toprule
case & AC-OPF & O-DC-OTS & AC-OTS & ours & best known & bound & gap \\
\midrule
\endhead
\code{3-lmbd}        & \num{11236}   & \num{10636}   & \num{10636}   & \num{10636}   & \num{10636}   & \num{10636}   & $0.00\%$ \\
\code{5-pjm}         & \num{76377}   & \num{75190}   & \num{75190}   & \num{75190}   & \num{75190}   & \num{73376}   & $2.41\%$ \\
\code{14-ieee}       & \num{5999}    & \num{5999}    & \num{5999}    & \num{5999}    & \num{5999}    & \num{5735}    & $4.39\%$ \\
\code{24-ieee-rts}   & \num{134944}  & \num{122283}  & \num{119743}  & \num{119743}  & \num{119743}  & \num{111741}  & $6.68\%$ \\
\code{30-as}         & \num{4996}    & \num{2797}    & \num{2797}    & \num{2797}    & \num{2797}    & \num{2767}    & $1.06\%$ \\
\code{30-ieee}       & \num{18044}   & \num{17939}   & \num{17936}   & \num{17940}   & \num{17936}   & \num{17287}   & $3.62\%$ \\
\code{39-epri}       & \num{249672}  & \num{246850}  & \num{246723}  & \num{246561}  & \num{246561}  & \num{245369}  & $0.48\%$ \\
\code{57-ieee}       & \num{49290}   & \num{49290}   & \num{49274}   & \num{49272}   & \num{49272}   & \num{49224}   & $0.10\%$ \\
\code{60-c}          & \num{185239}  & \num{182028}  & \num{182028}  & \num{182028}  & \num{182028}  & \num{180836}  & $0.65\%$ \\
\code{73-ieee-rts}   & \num{422627}  & \num{413133}  & \num{385194}  & \num{389142}  & \num{385194}  & \num{368338}  & $4.38\%$ \\
\code{89-pegase}$^\dagger$  & \num{130175} & \num{100702} & \num{100344} & \num{100503} & \num{100344} & \num{99858} & $0.48\%$ \\
\code{118-ieee}      & \num{242237}  & \num{195918}  & \num{180312}  & \num{222376}  & \num{180312}  & \num{169523}  & $5.98\%$ \\
\code{179-goc}$^\dagger$    & \num{1932044} & \num{1931004} & --        & \textbf{\num{1928125}} & \num{1928125} & \num{1835909} & $4.78\%$ \\
\code{200-activ}     & \num{35701}   & \num{35701}   & \num{35701}   & \num{35701}   & \num{35701}   & \num{35670}   & $0.09\%$ \\
\code{240-pserc}$^\dagger$ & \num{4640589} & \num{4627155} & --      & \num{4635816} & \num{4627155} & \num{4606386} & $0.45\%$ \\
\code{300-ieee}$^\dagger$   & \num{684985} & \num{684985} & \num{683968} & \num{684401} & \num{683968} & \num{678283} & $0.83\%$ \\
\code{500-goc}       & \num{692407}  & \num{692271}  & --            & \textbf{\num{676183}}  & \num{676183}  & \num{667141}  & $1.34\%$ \\
\code{1354-pegase}$^\dagger$ & \num{1498271} & \num{1496750} & -- & \num{1498117} & \num{1496750} & \num{1487079} & $0.65\%$ \\
\bottomrule
\end{longtable}

\noindent
$^\dagger$ Rigorous bound on the exact (unmerged) model; see Section~\ref{sec:rigour}. All other rows
never involved merging, so \emph{every} bound in this table is rigorous.

\paragraph{Solutions improving on the reference.} \code{179-goc} \num{1928125} and \code{500-goc}
\num{676183} (both new; Section~\ref{sec:cardinality}), \code{39-epri} \num{246561} (opening lines $\{4,6\}$,
better than the paper's Juniper AC-OTS \num{246723}); \code{500-goc} \num{685131} (better than the
paper's O-DC-OTS \num{692271}); \code{57-ieee} \num{49272}; \code{300-ieee} \num{684401}.

\paragraph{Where the relaxation is weak.} The largest gaps are on \emph{small}, highly combinatorial
cases (\code{24-ieee-rts} \SI{6.68}{\percent}, \code{118-ieee} \SI{5.98}{\percent}, \code{179-goc}
\SI{4.92}{\percent}, \code{14-ieee} \SI{4.39}{\percent}). Small does not mean easy here: a larger
fraction of the network is switchable relative to its redundancy.

\section{Certification matters as much as the relaxation}

Re-certifying every stored relaxation with the optimized certificate tightened every case, in several
substantially: \code{179-goc} \SI{6.28}{\percent} $\to$ \SI{4.88}{\percent}, \code{60-c}
\SI{1.08}{\percent} $\to$ \SI{0.65}{\percent}, \code{200-activ} \SI{0.30}{\percent} $\to$
\SI{0.09}{\percent}. Two conclusions:

\begin{itemize}
  \item \textbf{The big-M versus no-big-M ranking was a certification artifact.} Big-M appeared inferior
        because its certificates were looser, not because its relaxation was weaker. Certified properly,
        big-M gives the \emph{better} bound on \code{179-goc} (\num{1836683} versus \num{1810514}).
        No-big-M remains 1.5--2$\times$ faster to solve.
  \item \textbf{Tight trace bounds never helped}, because the residual term dominates the box
        certificate. A clean negative result.
\end{itemize}

After re-certification the certified bound lies within \SI{0.25}{\percent} of the raw SDP objective on
every relaxation, and within \SI{0.03}{\percent} on most, so the remaining gaps are attributable to the
relaxations rather than to inexact solves.

\section{Rigour of the bounds}\label{sec:rigour}

\subsection{A preprocessing step invalidated the bounds}

Low-impedance bus merging (\code{merge\_zmax}, default $10^{-3}$) was applied on 5 of the 18 cases. It is
\emph{not} a relaxation, for three independent reasons:

\begin{enumerate}
  \item \textbf{Restriction (voltages).} Buses joined by a merged tie are forced to share one $(e,f)$
        pair; the original problem does not require this.
  \item \textbf{Restriction (switching).} Merged ties are removed from the switchable set, i.e.\ forced
        closed. Minimizing over a \emph{subset} of switching configurations can only raise the optimum.
        This affects 178 branches on \code{1354-pegase}.
  \item \textbf{Relaxation (losses).} Merged ties are modelled as lossless, which acts in the opposite
        direction.
\end{enumerate}

Because the effects have opposite signs, the merged optimum is neither provably above nor below the true
one, so a bound computed on a merged model is not a valid lower bound for the original problem.

\subsection{The failure this caused}

On \code{1354-pegase} the consequence was not subtle. Merging at the default threshold merged a tie
sitting exactly at its thermal limit (flow/rate $=1.000$), which rendered the model \emph{infeasible}:
with every switch closed the merged POP is \code{LOCALLY\_INFEASIBLE}, while the unmerged one solves to
\num{1498271.03}, exactly the paper's AC-OPF cost. An SDP relaxation of an infeasible problem can take
any value, and indeed the re-certification reported ``lower bounds'' of \num{2580996}, \num{2612454} and
\num{3772084} against a known feasible cost of \num{1496750} --- exceeding it by 72\% to 152\%.

The mechanism was isolated by scaling only the merged ties' thermal ratings: at $\times 1.00$ and
$\times 1.02$ the model stays infeasible, at $\times 1.10$ it solves. Voltage differences across merged
ties are negligible at the true solution (maximum \num{0.00194}~p.u.), so the shared-voltage restriction
was nearly harmless; the binding thermal limit was not.

\paragraph{Remedies now in place.} A self-validating merge rule
(\code{safe\_merge\_exclusions}) compares the all-closed AC-OPF of the merged model against the unmerged
one and un-merges the most heavily loaded tie until they agree, terminating in the worst case with
nothing merged. On \code{1354-pegase} it un-merges 6 of 184 ties and reproduces the reference AC-OPF to
\SI{-0.018}{\percent}; on every other case it excludes nothing. A regression test
(\code{scripts/validate\_merging.jl}) checks all 18 cases, and the experiment scripts now verify that no
certified bound exceeds a known feasible cost, with the tabulation step failing outright if one does.

\subsection{Removing merging costs almost nothing}

Since the unmerged model \emph{is} the exact problem, and the certificate is valid at any dual point,
rigour is obtained simply by not merging. Re-running all variants on the exact model:

\begin{table}[h]\centering
\begin{tabular}{lrrrl}
\toprule
case & merged ties & merged bound & rigorous bound & certified gap \\
\midrule
\code{179-goc}   & 2  & \num{1836682.93} & \num{1835908.57} & $4.883\% \to 4.923\%$ \\
\code{89-pegase} & 19 & \num{99890.19}   & \num{99858.29}   & $0.452\% \to 0.484\%$ \\
\code{300-ieee}  & 2  & \num{678466.75}  & \num{678283.45}  & $0.804\% \to 0.831\%$ \\
\code{240-pserc} & 55 & \num{4606997.21} & \num{4606386.34} & $0.436\% \to 0.449\%$ \\
\code{1354-pegase} & 178 & \num{1486086.12} & \num{1487078.55} & $0.712\% \to \mathbf{0.646\%}$ \\
\bottomrule
\end{tabular}
\end{table}

\noindent
\textbf{Rigour is free.} Across the five affected cases the change in certified gap ranges from
$+0.040$ to $-0.066$ percentage points: it costs almost nothing on four cases and \emph{improves} the
largest, \code{1354-pegase}, where the exact model gives a strictly better bound than the merged one.
\textbf{All eighteen cases now carry rigorous bounds} (thirteen never merged anything; five re-run on
the exact model).

Two caveats. First, part of even that small deficit is budget rather than modelling: the merged numbers
come from runs with a \SI{1800}{\second} certificate budget and the rigorous ones from
\SI{900}{\second}; in a controlled comparison at equal budget the unmerged bound was \emph{better} on
\code{179-goc} and \code{300-ieee} (and five times faster on the latter). Second, merging is not
obsolete: the \code{89-pegase} big-M variant fails numerically when unmerged, returning a garbage point
and certificates near $-1.6\times 10^{8}$. That is still a valid (if useless) bound, and the no-big-M
variants supply the best bound anyway, but \code{merge\_zmax} cannot simply be set to zero and the merge
machinery deleted.

\section{Rounding at scale}

On small enumerable instances (Experiments 1--2, with ground truth by enumeration) correlated schemes
outperform independent rounding, confirming that order-2 pseudo-moments carry usable joint information.

At scale the picture changes sharply. Over 100 samples per scheme, the AC-feasible fraction is
39--\SI{55}{\percent} on \code{89-pegase} and 24--\SI{37}{\percent} on \code{240-pserc}, but
0--\SI{2}{\percent} on \code{118-ieee}, \code{179-goc}, \code{300-ieee}, \code{500-goc} and
\code{1354-pegase} --- under \emph{every} scheme, including island-guarded sampling and island repair.

\paragraph{The cause is cardinality, not connectivity.} Independent and Gaussian rounding decide each
line separately, so the number of lines opened simultaneously concentrates around $\sum_\ell (1 - y_\ell)$,
which is 25 on \code{300-ieee} and \num{724} of \num{1807} binaries on \code{1354-pegase}. Capping it,
while still letting the relaxation choose \emph{which} lines (weighted sampling without replacement by
$1-y_\ell$, implemented as a Gumbel top-$k$ draw):

\begin{table}[h]\centering
\begin{tabular}{lrrrrrl}
\toprule
case & independent & cap 8 & cap 5 & cap 3 & cap 1 & best cost found \\
\midrule
\code{118-ieee} & $0/50$ & \SI{0}{\percent}  & \SI{6}{\percent}  & \SI{22}{\percent} & \SI{58}{\percent} & \num{225486} ($+25.05\%$) \\
\code{300-ieee} & $0/30$ & \SI{47}{\percent} & \SI{67}{\percent} & \SI{77}{\percent} & \SI{93}{\percent} & \num{684766} ($+0.12\%$) \\
\bottomrule
\end{tabular}
\end{table}

On \code{300-ieee} this is the first time rounding alone, with no local-search polishing, has produced
competitive solutions at that size. On \code{118-ieee} feasibility is restored but solution quality is
not, and the reason is diagnostic: the relaxation implies roughly 11 openings whereas the best known
solution opens \num{37}, so \emph{no} sampler drawing from those fixed marginals can reach it. At that
size the limitation is the pseudo-distribution, not the rounding scheme.

\paragraph{Islanding is a partial explanation only.} Island repair does not restore feasibility, but the
dominant rejection reason varies by case. Of 100 independent samples, AC local infeasibility versus
island-screen rejection splits 78/22 on \code{118-ieee}, 78/18 on \code{200-activ} and 95/5 on
\code{500-goc}, but 18/82 on \code{179-goc}, 46/54 on \code{300-ieee} and 0/100 on \code{1354-pegase}.
Cardinality capping helps because it addresses both at once: opening fewer lines produces fewer islands
\emph{and} less network stress.

\subsection{How many lines should a sample open?}\label{sec:cardinality}

\code{scripts/experiment5\_cardinality.jl} sweeps the cap directly, reusing the stored marginals so no SDP
is re-solved. $N=100$ ($N=60$ on \code{500-goc}), evaluated on the original network.

\begin{table}[h]\centering
\begin{tabular}{lrrlrrr}
\toprule
case & E[\#opened] & independent & best cap & feasible & best cost & vs best known \\
\midrule
\code{118-ieee}  & 10.9 & \SI{0}{\percent}  & $k{=}3$ & \SI{17}{\percent} & \num{225486.25}  & $+25.05\%$ \\
\code{179-goc}   & 39.3 & \SI{0}{\percent}  & $k{=}2$ & \SI{78}{\percent} & \num{1928124.99} & $\mathbf{-0.15\%}$ \\
\code{200-activ} & 19.2 & \SI{5}{\percent}  & $k{=}2$ & \SI{77}{\percent} & \num{35700.87}   & $+0.00\%$ \\
\code{240-pserc} & 33.0 & \SI{26}{\percent} & $k{=}5$ & \SI{89}{\percent} & \num{4638056.84} & $+0.24\%$ \\
\code{300-ieee}  & 25.0 & \SI{0}{\percent}  & $k{=}1$ & \SI{91}{\percent} & \num{684584.26}  & $+0.09\%$ \\
\code{500-goc}   & 50.5 & \SI{0}{\percent}  & $k{=}6$ & \SI{52}{\percent} & \num{676182.83}  & $\mathbf{-1.31\%}$ \\
\code{89-pegase} & 37.2 & \SI{38}{\percent} & $k{=}8$ & \SI{74}{\percent} & \num{101016.95}  & $+0.67\%$ \\
\code{1354-pegase}$^{*}$ & 204.0 & \SI{0}{\percent} & $k{=}3$ & \SI{80}{\percent} & \num{1497390.91} & $+0.04\%$ \\
\bottomrule
\end{tabular}
\end{table}

\noindent
\textbf{Two new best-known solutions}, both verified independently by re-solving the AC-OPF for that
switching configuration from scratch on the original network:
\begin{itemize}
  \item \code{179-goc}: \num{1928124.99}, opening only two lines $\{147, 158\}$ --- \num{2851}
        (\SI{0.148}{\percent}) below our previous incumbent, and below the paper's O-DC-OTS
        (\num{1931004}) and AC-OPF (\num{1932044}).
  \item \code{500-goc}: \num{676182.83}, opening six lines $\{131, 228, 464, 607, 629, 698\}$ ---
        \num{8949} (\SI{1.306}{\percent}) below our incumbent and \SI{2.3}{\percent} below the paper's
        O-DC-OTS (\num{692271}).
\end{itemize}
Both tighten their certified gaps, since the best known cost falls: \code{179-goc}
$4.923\% \to 4.783\%$ and \code{500-goc} $2.626\% \to 1.337\%$, roughly halving the latter.

The effect is large and consistent: independent rounding lands at 0--\SI{38}{\percent} AC-feasible,
capping at 52--\SI{97}{\percent}. The useful caps are $1$--$8$, always far below the relaxation's own
expectation, which remains useless everywhere (0--\SI{2}{\percent}), and quality is flat across
neighbouring caps, so the result does not depend on tuning the cap. \code{118-ieee} stays the exception:
feasibility is restored but cost remains $+25\%$, because its optimum opens \num{37} lines while the
marginals imply about 11.

$^{*}$ \code{1354-pegase} needed fresh marginals: its original ones came from the pre-fix merged model,
which was infeasible (Section~\ref{sec:rigour}), and a dimension guard in the script caught this. They
were regenerated on the \emph{exact} unmerged model (experiment 4, \code{base}, \num{1991} binaries,
\SI{3.5}{\hour}). The difference is instructive --- the infeasible model implied
$\mathbb{E}[\#\text{opened}] = 724$ of \num{1807} binaries, the exact model $204$ of \num{1991} --- but it
means this row rests on exact-model marginals while the others use merged-model ones.

With \code{1354-pegase} included the finding is unanimous across all eight cases: independent rounding is
hopeless (0--\SI{38}{\percent}), capping restores feasibility (52--\SI{97}{\percent}), and the useful cap
is $1$--$8$ regardless of network size.

\section{Negative results worth recording}

\begin{itemize}
  \item Tight trace bounds in the box certificate never changed the answer.
  \item Island repair and guarded sampling help only where rounding already works.
  \item Connectivity cuts and radial fixings improve the bound on \code{89-pegase} and
        \code{1354-pegase} but do not improve rounding on \code{118-ieee}, \code{179-goc} or
        \code{300-ieee}.
  \item Capping openings at the relaxation's own expectation is insufficient (\SI{2}{\percent} and
        \SI{0}{\percent}); the useful caps are far tighter, 1--8 lines.
  \item The ``api'' cases are genuinely tight: on \code{118-ieee}, 100 of the 186 single-line openings
        are AC-infeasible on their own, and feasibility is \emph{not} monotone in the set of opened
        lines.
\end{itemize}

\section{Open problems}

\begin{enumerate}
  \item Sweep the opening cap $k_{\max}$ and re-run Experiments 3 and 4 with cardinality-limited
        sampling.
  \item For cases requiring large coordinated rewirings (\code{118-ieee}, \code{89-pegase},
        \code{73-ieee-rts}), move from resampling fixed marginals to conditioning and re-solving
        (randomized diving), which changes the distribution; and test conditional moments
        $\mathbb{E}[V \mid z]$ as Ipopt warm starts.
  \item Understand why the \code{118-ieee} marginals are so concentrated (expected 11 openings against
        an optimum opening 37). This is a relaxation-tightness question and connects to the
        relaxation-selection thread.
  \item The ML-guided relaxation-selection thread remains untouched; non-ML baselines (OBBT, QC
        envelopes, dual- and eigenvalue-guided order elevation) should come first.
\end{enumerate}

\section*{Artifacts}

Julia package \code{MomentMINLP} (JuMP, MOSEK, PowerModels, Ipopt/HSL). Supporting notes:
\code{results/writeup.md} (the Markdown counterpart of this document),
\code{results/merging\_findings.md}, \code{results/rounding\_scale\_findings.md},
\code{results/experiment4\_findings.md}, \code{results/recertify\_findings.md},
\code{results/numerics\_findings.md}. Regression tests: \code{scripts/validate\_merging.jl} (18/18) and
\code{scripts/test\_islands.jl} (96/96).

\begin{thebibliography}{9}
\bibitem{oustry2022}
A.~Oustry, C.~D'Ambrosio, L.~Liberti and M.~Ruiz,
\newblock ``Certified and accurate SDP bounds for the ACOPF problem,''
\newblock \emph{Power Systems Computation Conference (PSCC)}, 2022.
\end{thebibliography}

\end{document}
